\subsubsection{CONSISTÊNCIA EVENTUAL FORTE (\en{SEC})}

Consistência eventual forte define a consistência de réplicas sob o mesmo conjunto de atualizações, mesmo que essas atualizações sejam aplicadas em ordens distintas, desde que a casualidade seja respeitada. Interpretação de operações é, portanto, a base para a compreenção semâtica das operações entre as réplicas, pois garante coerência de contexto causal entre as réplicas. \en{SEC} assume que a interpretação de operações é correta, isso implica que as semâticas de domínio (especificações de cada sistema) devem ser preservadas para garantir a convergência \cite{Gomes2017VerifyingSEC}.

Gomes et al. (2017) apresentam exemplos de como a convergência de estado pode ser alcaçada por meio de regras que preservam a semâtica das operações em cenários concorrentes. Um desses exemplos é entitulado \en{Observable-Remove Set} (\en{OR Set}), este define remoções de elementos observáveis, isso é: cada elemento pode ser caracterizado por um identificador, a operação de remoção atua apenas sobre elementos nos quais o identificador é conhecido no momento de remoção; caso esses elementos removidos sejam adicionados novamente por uma operação concorrente, essa adição não será eliminada devido a discrepância de identificadores \cite{Gomes2017VerifyingSEC}. Esta implementação demonstra como a sincronização pode ser alcançada conforme as semânticas de domínio de cada sistema.

Ainda nos exemplos de Gomes et al. (2017), é descrito mecânismo de edição de texto que permite manter a caracteres excluídos, assim réplicas podem usar qualquer caractere como referência para inserir antes ou depois sem precisar considerar ações concorrentes. Isso pode ser aplicado na edição de anotações da aplicação móvel.

Nesse contexto, \en{CRDTs} garantem consistência eventual forte por especificar comutação de operações concorrentes, assim as réplicas podem convergir sem a necessidade de coordenação entre elas \cite{Shapiro2015CRDT}. 
